\section{Statutory Recommendation: Total Population with Prison Adjustment}
\label{sec:recommendation}

\subsection{The Recommendation}

Based on the empirical findings of this paper, we recommend \textbf{total population with prison adjustment only} as the preferred population base for congressional redistricting. Specifically, we recommend:

\begin{enumerate}
    \item Use 2020 Census total population (P.L. 94-171, Table P1) as the starting point
    \item Apply the prison-gerrymandering correction: reallocate Group Quarters types 101--103 (federal and state prisons, local jails) from facility tracts to estimated home-county tracts using the methodology described in N.1~\citep{dellalibera2026n1}
    \item Do not apply the college-student adjustment (campus counting remains the default)
    \item Do not substitute citizen VAP for total population
\end{enumerate}

\subsection{Rationale}

\textbf{Legal defensibility.} Total population is the constitutional default under \emph{Reynolds v. Sims} and has been used for congressional redistricting throughout American history. The prison-count adjustment is a legitimate correction to a documented counting distortion: incarcerated individuals are counted where they cannot vote, cannot access local services as residents, and have no attachment to the facility county. New York, California, and a majority of states have adopted prison-count corrections for state redistricting. The Census Bureau has announced plans to produce a prison-adjusted redistricting data file for 2030~\citep{census2022prison_adjustment}, providing federal endorsement for the correction methodology.

\textbf{Small redistricting disruption.} The prison adjustment produces only 0.8\% mean tract reassignment nationally---the second-smallest of the four comparisons. The disruption is concentrated in prison-dense rural counties, not broadly distributed across all districts. The adjustment corrects a real distortion (rural inflation) without creating new distortions.

\textbf{Compactness neutrality.} RSI changes under prison adjustment are below 0.5\% in all 43 redistricting states. Prison adjustment does not require districts to expand into less-compact territory.

\textbf{Data quality.} Total population from the decennial census is the highest-quality population data available: a near-complete enumeration with near-zero sampling error at the tract level. The prison adjustment uses Group Quarters data from the same decennial census, avoiding the ACS sampling error that plagues CVAP-based redistricting.

\subsection{Why Not Citizen VAP?}

\textbf{Systematic partisan distortion.} Citizen VAP produces a $-1.2$pp mean partisan direction shift toward Republican districts among affected districts, reaching $-3.1$pp in California. This is the largest partisan distortion of the four alternatives considered. It is consistent, predictable, and concentrated in the states where congressional redistricting decisions are most consequential.

\textbf{Data quality problems.} CVAP data derives from the ACS 5-year estimates, which carry substantial sampling error at the tract level. In 15--20\% of tracts in high-immigration states, the CVAP margin of error exceeds 15\% of the point estimate. Congressional redistricting at $\pm$0.5\% balance tolerance cannot be reliably achieved with population data of this quality.

\textbf{Constitutional exposure for congressional maps.} As argued in N.3~\citep{dellalibera2026n3}, citizen-VAP congressional redistricting in high-immigration states would produce total-population disparities across districts exceeding 100,000 persons in some cases---likely unconstitutional under \emph{Wesberry v. Sanders} and \emph{Karcher v. Daggett}.

\textbf{Large redistricting disruption.} The 3.2\% mean national reassignment rate, reaching 11.8\% in California, is four times larger than the prison adjustment's disruption. Adoption of citizen-VAP redistricting would require redrawing the boundaries of approximately 65,000 census tracts nationally relative to the total-population baseline.

\subsection{Why Not College Adjustment?}

The college-student adjustment is not recommended as part of the primary standard because:
\begin{enumerate}
    \item The case for counting students at home is weaker than the case for counting prisoners at home (students chose campus; prisoners did not)
    \item The data for home-state redistribution (IPEDS state-of-origin enrollment) is available only at the state level, not the county or tract level, introducing imprecision that the prison adjustment avoids
    \item College-student adjustment has not been adopted by any state legislature for redistricting, whereas prison-count reform has majority-state adoption
    \item The combined prison-plus-college adjustment produces a directionally mixed effect (prison correction is Democratic-direction; college correction is Republican-direction) that complicates the policy narrative
\end{enumerate}

The college-student adjustment remains available in BISECT as a sensitivity analysis option and may be appropriate for academic research or legislative exploration, but it is not recommended as a default redistricting standard.

\subsection{Alignment with Census Bureau 2030 Plans}

The Census Bureau has announced that for the 2030 redistricting data cycle, it will produce a prison-adjusted redistricting file that reallocates incarcerated populations to their home addresses~\citep{census2022prison_adjustment}. This planned product aligns precisely with our recommended approach: total population with prison adjustment. Adopting this recommendation now positions redistricting commissions and legislatures to transition smoothly to the 2030 Census Bureau file without methodology disruption.
